\documentclass{article}
\usepackage[utf8]{inputenc}
\usepackage[russian]{babel}
\usepackage{amsmath}
\usepackage{amsfonts}

\title{ADR-0001: Плоская архитектура и шаблон ReaderT Env IO}
\author{Оперативная группа проекта SMP}
\date{Июль 2026}

\begin{document}

\maketitle

\section*{Статус}
\textbf{ACCEPTED} (Принято к исполнению)

\section*{Контекст}
В процессе проектирования L7-балансировщика трафика \texttt{stm-markov-proxy} возникла необходимость координации статических настроек (координаты прокси, пулы бэкендов) и динамического многопоточного состояния. 

Традиционный для Haskell подход с многоэтажными трансформерами монад (например, стеки с тяжелым \texttt{StateT}) на практике приводит к катастрофическому оверхеду при масштабировании, усложняет сигнатуры функций и не является потокобезопасным из коробки в многопоточной среде GHC Runtime.

\section*{Решение}
В соответствии с методологией декомпозиции практических задач на Haskell Ивана Бабкина, утверждается плоская архитектура приложения:
\begin{enumerate}
    \item Создается единая монолитная структура данных \texttt{Env} (Регистр), инкапсулирующая как статический конфиг, так и изменяемые потокобезопасные ссылки.
    \item В качестве базового монадического стека приложения утверждается единственный плоский трансформер:
    \[ \text{type } ProxyM\ a = ReaderT\ Env\ IO\ a \]
    \item Доступ к контексту осуществляется через иммутабельный вызов \texttt{ask}, а изменение динамических параметров делегируется на уровень транзакционной памяти.
\end{enumerate}

\section*{Последствия}
\begin{itemize}
    \item \textbf{Плюсы:} Полное стирание монадической обертки компилятором GHC при оптимизации, предсказуемый проброс указателя на \texttt{Env} через регистры процессора, отсутствие бойлерплейта в сигнатурах.
    \item \textbf{Минусы:} Необходимость явного вызова функций лифтинга (\texttt{liftIO}) для выполнения эффектов реального мира на уровне обработчиков.
\end{itemize}

\end{document}
